\section{Conclusions}
\label{sec:conclusions}
SDALR's high same-bearing accuracy does not carry over to unseen physical bearings under the tested protocol. On Paderborn, the median loss is 36.2 points across ten pre-specified splits, and 36.6 points for SHOT on the same splits; in a post-hoc within-bearing contrast, a given bearing loses a median 38.7 points when unseen rather than seen. Non-adaptive forests lose 22.7--32.1 points, a same-condition forest control 29.7, and SDALR on HUST 33.5, although HUST changes bearing size with specimen. These effects describe the available bearing pools rather than a general population guarantee.

A post-hoc source-only analysis finds a 35.1-point gap already in the source model, and adaptation does not eliminate it. Expanding a forest's source from one to four specimens per class raises unseen-bearing accuracy from 46 to 75\,\%, while training volume grows simultaneously. This supports the value of a larger source collection but does not distinguish insufficient diversity from sample count or reliance on specimen-associated features. The kinematic forest's median gap is 8.4 points in ordinary and 10.4 points in balanced accuracy, on either side of its pre-specified 10-point threshold, so it does not establish transfer.

Envelope gates falsely flag 0 to 9 of 480 recordings from six healthy bearings, against 214 for the tested simplified raw-spectrum rule. The envelope threshold gate adds a median 4.5 points inside SDALR, closing about a fifth of the gap recovered by an oracle filter. In a post-hoc comparison, used directly on the recordings it certifies, its verdicts add 8.7 points, more than inside adaptation in every split. Most gain occurs on certified recordings, and coverage is concentrated on larger-damage specimens under this pipeline. These observations establish neither general safety nor an intrinsic detectability limit for smaller faults.

A credible unseen-bearing claim should therefore include bearing-disjoint targets across source splits, source-only and adapted results, comparable class weighting, paired module differences, a precisely defined oracle, shallow baselines, false acceptance and coverage by specimen, and transparent implementation and model-selection checks. The indexing defect found and corrected in our own SHOT implementation illustrates the need for those checks; the SFDA conclusions are limited to SDALR and SHOT.
